% =============================================================================
%  RL-H2-02 · Tabel positioning vs paper terdekat
% -----------------------------------------------------------------------------
%  Role    : RL — Research Lead (@alim)
%  Sprint  : H2 (setelah RL-H1-02) · PARALEL
%  Ticket  : RL-H2-02                        Status: DONE · 2026-11-24
%  Target  : docs/positioning.tex  ->  docs/positioning.pdf
%  Input   : docs/lit/*.tex + abstrak + brief v2 §2/§6
%  Output  : dokumentasi pembeda (11 paper pesaing + Studi ini)
%  Handoff : RL-H4-01 (related work), laporan akhir Sprint 1
%  DoD     : [x] baris "studi ini" hanya berisi klaim AMAN
%            [x] setiap sel pesaing bersumber (lit note/abstrak)
%            [x] 3 kalimat posisi konsisten dengan batas klaim brief §6
% =============================================================================
\documentclass[11pt,a4paper]{article}
% =============================================================================
%  00_preamble.tex | shared preamble untuk semua deliverable RL
%  Job     : bukan job card; dipakai oleh RL-H1-01 .. RL-H7-01
%  Target  : docs/ (salin bersama file .tex)
%  Catatan : kompilasi dengan pdflatex. Untuk XeLaTeX/LuaLaTeX, hapus inputenc.
% =============================================================================

\usepackage[utf8]{inputenc}
\usepackage[T1]{fontenc}
\usepackage{lmodern}
\usepackage[a4paper,margin=2.4cm]{geometry}
\usepackage{amsmath,amssymb}
\usepackage{booktabs}
\usepackage{longtable}
\usepackage{array}
\usepackage{graphicx}
\usepackage{xcolor}
\usepackage{enumitem}
\usepackage{float}
\usepackage{caption}
\usepackage{url}
\usepackage{microtype}
\usepackage[colorlinks=true,linkcolor=blue!50!black,urlcolor=blue!50!black,citecolor=blue!50!black]{hyperref}

% --- helper wajib -----------------------------------------------------------
% Path/nama file aman (underscore, garis miring) di dalam teks.
\newcommand{\fpath}[1]{\texttt{\detokenize{#1}}}
% TODO yang mencolok. Jangan dihapus sebelum job selesai.
\newcommand{\todo}[1]{\textcolor{red!75!black}{\textbf{[TODO: #1]}}}
\newcommand{\jobid}[1]{\textsf{\textbf{#1}}}
\newcommand{\runid}[1]{\texttt{\detokenize{#1}}}
\newcommand{\claim}[1]{\textsf{#1}}
\newcommand{\good}{\textcolor{green!45!black}{$\checkmark$}}
\newcommand{\bad}{\textcolor{red!75!black}{$\times$}}
\newcommand{\maybe}{\textcolor{orange!80!black}{$\sim$}}

\newcolumntype{L}[1]{>{\raggedright\arraybackslash}p{#1}}
\newcolumntype{C}[1]{>{\centering\arraybackslash}p{#1}}

\setlength{\parindent}{0pt}
\setlength{\parskip}{4pt}
\captionsetup{font=small,labelfont=bf}
\setlist[itemize]{leftmargin=1.2em,itemsep=2pt,topsep=2pt}
\setlist[enumerate]{leftmargin=1.5em,itemsep=2pt,topsep=2pt}

% Judul job yang konsisten di semua dokumen.
\newcommand{\jobheader}[4]{%
  \begin{center}
  {\Large\bfseries #1}\\[2pt]
  {\large #2}\\[4pt]
  \begin{tabular}{@{}ll@{}}
  \toprule
  Role & #3 \\
  Sprint & #4 \\
  \bottomrule
  \end{tabular}
  \end{center}
}


\title{\jobid{RL-H2-02} | Positioning vs 11 Paper Terdekat \& Strategi Batas Klaim\\
\large Apa yang Sudah Ada, Apa yang Belum, dan Peluang Kontribusi Studi Ini}
\author{RL (Alim) \\ \small RENARO Lab}
\date{24 November 2026 (H2) \quad | \quad \textcolor{green!45!black}{STATUS: FINAL H2}}

\begin{document}
\maketitle

\section{Properti graph (11 paper terdekat vs studi ini)}

\begin{table}[H]
\centering
\footnotesize
\setlength{\tabcolsep}{3pt}
\begin{tabular}{@{}L{2.3cm}L{3.9cm}C{1.4cm}C{1.4cm}L{2.5cm}L{2.9cm}@{}}
\toprule
Paper & Domain / data & Berarah & Dinamis & Jeda per edge ($\ell$) & Kondisi pembentuk graph \\
\midrule
RCSGNet \cite{bai2026rcsgnet} & Kentang, 25 provinsi Tiongkok, mingguan 2012--2018
& ya & ya & tidak (jendela saat ini) & laten (state pasar $s_t$ dari harga) \\
MarketDNA \cite{marketdna2026} & Kentang multi-centre, harian \& periodik
& ya & ya & tidak & laten diskret (token state harga) \\
MillGNN \cite{millgnn2025} & Deret waktu multi-skala \emph{lead-lag} umum
& ya & ya & ya (diskret + bias $-\lambda\ell$) & teramati dari jendela input (tanpa iklim) \\
SAMSGL \cite{samsgl2023} & Spatio-temporal umum dgn \emph{time delay}
& ya & ya & ya (\emph{series-aligned} FFT) & teramati dari jendela input (tanpa iklim) \\
CRAFormer \cite{craformer2025} & Pergeseran tanah longsor, GNSS + hujan
& ya (DAG) & semi-statis & ya (per-lag DAG fisik) & agregat tetap / mask 5 peran kausal \\
DRAN \cite{dran2025} & Peramalan spatio-temporal umum
& ya & ya & tidak & gate relasi statis + dinamis \\
SDGF \cite{sdgf2025} & Deret waktu multivariat multi-skala
& ya & ya & tidak & \emph{attention gate} statis + dinamis \\
STEVE \cite{ji2025steve} & Trafik perkotaan (\emph{robust traffic})
& ya & ya & tidak & \emph{confounder} laten dari histori \\
CGNN \cite{cgnn2023} & Harga pertanian India (kentang \& tomat)
& tidak & tidak & tidak & tidak ada (graph jarak statis) \\
Min et al.\ (2025) & Harga grosir pertanian Korea Selatan
& ya & ya & tidak (\emph{smoothing window}) & laten (\emph{self-attention} StemGNN) \\
Teh Sri Lanka \cite{tea2026} & Harga teh Sri Lanka, lag cuaca sumber
& tanpa graph & tanpa graph & lag fitur 1--3 minggu & tanpa graph (regresi deret waktu) \\
\midrule
\textbf{Studi ini} & \textbf{Cabai merah keriting, 34 provinsi Indonesia, harian kerja (2018--2026)}
& \textbf{ya} & \textbf{ya} & \textbf{ya (distribusi $\ell \in \{0..8\}$ pentad + \texttt{null} per edge)} & \textbf{teramati: iklim fisik di provinsi asal ($\mathbf{c}_{j,t}$)} \\
\bottomrule
\end{tabular}
\caption{Perbandingan properti graph terhadap 11 studi terdekat. Baris ``Studi ini'' hanya memuat klaim Kategori AMAN (Brief v2 \S6).}
\end{table}

\section{Properti evaluasi \& ketahanan (\emph{robustness})}

\begin{table}[H]
\centering
\footnotesize
\setlength{\tabcolsep}{4pt}
\begin{tabular}{@{}L{2.6cm}L{3.2cm}L{3.8cm}L{4.4cm}@{}}
\toprule
Paper & Iklim eksogen & Uji OOD regime iklim & Uji graph prior keliru \\
\midrule
RCSGNet \cite{bai2026rcsgnet} & tidak (\emph{price-only}) & tidak (split kronologis standar) & tidak \\
MarketDNA \cite{marketdna2026} & tidak (\emph{price-only conditions}) & tidak & tidak \\
MillGNN \cite{millgnn2025} & tidak & tidak & tidak \\
SAMSGL \cite{samsgl2023} & tidak (pada struktur graph harga) & tidak & tidak \\
CRAFormer \cite{craformer2025} & ya (hujan lokal longsor) & tidak (fokus titik badai lokal) & tidak \\
DRAN \cite{dran2025} & tidak & tidak (pergeseran distribusi umum) & tidak \\
SDGF \cite{sdgf2025} & tidak & tidak & tidak \\
STEVE \cite{ji2025steve} & tidak langsung (\emph{unseen confounder}) & ya, tetapi OOD trafik (bukan rezim ENSO/kekeringan) & tidak \\
CGNN \cite{cgnn2023} & ya (hanya fitur node lokal) & tidak & tidak \\
Min et al.\ (2025) & ya (hanya fitur node/variabel) & tidak & tidak \\
Teh Sri Lanka \cite{tea2026} & ya (lag cuaca wilayah sumber) & tidak & tidak (tanpa graph) \\
\midrule
\textbf{Studi ini} & \textbf{ya (iklim provinsi asal + lokal, \emph{as-of} \texttt{available\_at <= t})} & \textbf{ya (stratifikasi \texttt{dry}, \texttt{extreme\_wet}, \texttt{el\_nino}, \texttt{la\_nina}, \texttt{normal} + \emph{Worst-Regime MAE})} & \textbf{ya (perturbasi \texttt{drop/add/rewire} 10--50\% pada \texttt{bps\_mask} + prior \texttt{none})} \\
\bottomrule
\end{tabular}
\caption{Perbandingan properti evaluasi. Kolom OOD memisahkan uji rezim iklim fisik eksogen dari uji pergeseran waktu/trafik biasa.}
\end{table}

\section{Sumber setiap baris pesaing}
\begin{table}[H]
\centering
\footnotesize
\begin{tabular}{@{}L{2.8cm}L{5.8cm}L{5.6cm}@{}}
\toprule
Paper & Sumber Dokumen di Repo / Literatur & Halaman / Bagian yang Diverifikasi \\
\midrule
RCSGNet & \fpath{docs/lit/rcsgnet.tex} \& Dossier Bag.\ I & Abstrak verbatim \& spesifikasi 4 komponen (\emph{Potato Research} Vol.\ 69 No.\ 5, Juni 2026) \\
MarketDNA & \fpath{docs/lit/marketdna.tex} \& Dossier Bag.\ II & Abstrak verbatim \& spesifikasi 3 komponen (\emph{Potato Research}, Agustus 2026) \\
MillGNN & \fpath{docs/lit/millgnn.tex} (\texttt{arXiv:2509.03852}) & Abstrak \& formulasi \emph{decay-aware lead-lag attention} ($-\lambda\ell$) \\
SAMSGL & \fpath{docs/lit/millgnn.tex} (\texttt{arXiv:2312.02646}) & Abstrak \& \emph{series-aligned graph convolution} berbasis FFT \\
CRAFormer & \fpath{docs/references.bib} (\emph{Entropy} Vol.\ 28 No.\ 1) & Zhang dkk.\ (2025/2026): \emph{Dynamic-Lag Causal Graphs} untuk longsor dgn driver hujan \\
DRAN & \fpath{docs/lit/sdgf.tex} (\texttt{arXiv:2504.01531}) & Abstrak \& cabang adaptasi distribusi/relasi statis-dinamis \\
SDGF & \fpath{docs/lit/sdgf.tex} (\texttt{arXiv:2509.18135}) & Abstrak \& fusi korelasi statis dan dinamis multi-skala \\
STEVE & \fpath{docs/lit/steve.tex} (Prosiding \emph{KDD 2025}) & Ji dkk.\ (2025): \emph{basis confounder representations} pada prediksi trafik OOD \\
CGNN & \fpath{docs/references.bib} (\texttt{arXiv:2304.09761}) & Abstrak \& metodologi CNN-GNN harga kentang/tomat India dgn fitur iklim node \\
Min et al.\ 2025 & Brief v2 \S2 \& indeks PMC/NIH (2025) & Perbandingan StemGNN/TGCN vs LSTM harga pertanian Korea Selatan + cuaca + \emph{smoothing} \\
Teh Sri Lanka & \fpath{docs/references.bib} (\texttt{arXiv:2608.24894}) & Abstrak \& pemodelan lag cuaca wilayah sumber 1--3 minggu tanpa graph \\
\bottomrule
\end{tabular}
\end{table}

\section{Tiga kalimat posisi resmi (konsisten dengan Brief v2 \S6)}
\begin{enumerate}
  \item \textbf{Klaim Kontribusi Inti (Kategori AMAN)}: Sepengetahuan kami, belum ada model \emph{spatio-temporal graph neural network} (ST-GNN) antarprovinsi untuk peramalan harga komoditas pangan strategis di Indonesia yang secara eksplisit mengondisikan \emph{distribusi probabilitas jeda transmisi per edge berarah} ($\ell \in \{0..8\}$ pentad beserta slot tidak-aktif/\texttt{null}) pada \emph{indikator iklim fisik yang teramati di provinsi asal/pemasok} ($\mathbf{c}_{j,t}$: anomali hujan pentad CHIRPS v3, durasi hari kering beruntun, kejadian hujan ekstrem, dan indeks ENSO RONI).
  \item \textbf{Klaim Benchmark \& Ketahanan (Kategori AMAN)}: Studi ini menyediakan kerangka benchmark evaluasi harga pangan antarprovinsi yang bebas kebocoran waktu (\emph{availability-safe} melalui aturan \emph{as-of join} \texttt{available\_at <= t} yang terverifikasi secara empiris), dievaluasi secara terstratifikasi lintas rezim iklim fisik (\texttt{dry}, \texttt{extreme\_wet}, \texttt{el\_nino}, \texttt{la\_nina}, \texttt{normal}) menggunakan metrik \emph{Worst-Regime MAE}, serta diuji secara terkontrol terhadap kesalahan spesifikasi prior graph perdagangan (\emph{misspecified graph prior} melalui perturbasi struktural 10--50\% dan ablasi tanpa prior).
  \item \textbf{Pembatasan Klaim \& Interpretasi (Kategori JANGAN DIKLAIM)}: Kami \textbf{tidak} mengklaim pembelajaran graph berarah/adaptif, mekanisme \emph{gating} statis-dinamis, regularisasi invarian, ataupun penyertaan variabel cuaca ke dalam GNN harga pertanian sebagai hal baru; seluruh bobot edge dan distribusi jeda yang dipelajari diinterpretasikan sebagai \emph{dependensi prediktif berarah yang bergantung pada kondisi iklim} (\emph{climate-conditioned directional predictive dependency}), bukan sebagai arus pasokan fisik maupun bukti kausalitas struktural.
\end{enumerate}

\section{Batas klaim \& strategi membuka peluang temuan eksperimental}
\begin{table}[H]
\centering
\footnotesize
\begin{tabular}{@{}L{2.2cm}L{6.0cm}L{6.0cm}@{}}
\toprule
Kategori & Boleh diklaim (dengan kualifikasi yang tepat) & Tidak boleh diklaim \\
\midrule
\textbf{AMAN} & Selalu diawali \emph{``Sepengetahuan kami...''}: pengondisian distribusi jeda per edge ($\ell \in \{0..8\}$ pentad + slot \texttt{null}) pada iklim fisik teramati di provinsi asal ($\mathbf{c}_{j,t}$); protokol evaluasi \emph{availability-safe} terstratifikasi rezim iklim dgn \emph{Worst-Regime MAE}; serta uji terkontrol melawan baseline rezim laten ala RCSGNet (\jobid{MLE-H4-02}) dan lag-graph tanpa iklim ala MillGNN (\jobid{MLE-H4-03}) pada backbone yang sama. & Mengklaim kontribusi berlaku universal tanpa kualifikasi konteks jaringan harga pangan antarprovinsi Indonesia. \\
\midrule
\textbf{BERISIKO} (\emph{Peluang Upside Eksperimental}) & Melaporkan efektivitas cabang \emph{adaptive residual graph} + \emph{reliability gate} (Modul 2--3) serta objektif ketahanan lintas-rezim (Modul 4) sebagai \textbf{temuan eksperimental / empiris} dalam menahan degradasi MAE ketika prior graph BPS mengalami perturbasi 10--50\% (\jobid{CGS-H5-01}) atau ketika terjadi pergeseran rezim iklim ekstrem. & Mengklaim mekanisme \emph{adaptive graph}, \emph{static-dynamic gating}, atau \emph{invariant regularization} sebagai kontribusi metode/arsitektur baru. \\
\midrule
\textbf{JANGAN DIKLAIM} & --- & (1) ``Pertama memakai variabel iklim/cuaca pada GNN harga pertanian'' (sudah ada CGNN \cite{cgnn2023} \& Min et al.\ 2025); (2) ``Pertama memodelkan rezim pada GNN harga komoditas'' (sudah ada RCSGNet \cite{bai2026rcsgnet} \& MarketDNA \cite{marketdna2026}); (3) Mengklaim bobot edge $\mathbf{A}_t^{(\ell)}$ sebagai arus pasokan fisik tonase atau kausalitas intervensi. \\
\bottomrule
\end{tabular}
\end{table}

\end{document}
